%% kapa: exam=1 round=37 year=2026 period=1 form=A subject=real_estate subject_ko=부동산학원론 seq=40 number=120
%% source: Q-Net 감정평가사 기출문제 (cst00309) artlSeq=5258924 file=raw/1008/5258924/2026년 제37회 감정평가사 1차 1교시.pdf
%% conversion: kordoc 4.24.51 -> tools/gen1.py (auto); review=proofread-sonnet; flags=none
\begin{question}{120}
부동산신탁에 관한 설명으로 옳지 않은 것은?
\begin{choices}
  \item 신탁부동산은 수탁자 개인재산으로부터 독립한 재산으로 취급되므로 개인 채권자의 강제집행이 배제된다.
  \item 신탁기간 중 소유자가 자금이 필요하면 일정한 조건 하에 신탁수익권을 양도하거나 담보하여 자금을 차입할 수 있다.
  \item 관리신탁은 해외교포나 고령자 등 부동산을 직접 관리하기 곤란하거나 전문지식이 부족한 경우에 주로 활용된다.
  \item 처분신탁은 처분방법이나 절차에 어려움이 있는 부동산이나 대형ㆍ고가의 부동산을 효율적으로 처분할 수 있다.
  \item 담보신탁은 금융기관 입장에서 부동산관리가 자유로우나 법원 경매에 비해 환가절차가 복잡하다.
\end{choices}
\end{question}
\answer{5}
